\section{Moltiplicazione di polinomi e FFT - Fast Fourier Transform}
{
	\footnotesize Nota: la DFT (Discrete Fourier Transform) è la trasformata matematica discreta di Fourier, la FFT (Fast Fourier
	Transform) è l'algoritmo efficiente per calcolare la DFT.
}

\subsubsection*{Cenni sulla moltiplicazione di polinomi}
Dati due polinomi \(p(x)\) e \(q(x)\) rispettivamente di degree bound \(n\) e \(m\), il loro prodotto \(p(x) \cdot q(x)\)
è un polinomio di degree bound \(n + m\) il cui vettore di coefficienti \(c\) è dato dal prodotto di convoluzione tra i vettori
di coefficienti \(a\) e \(b\) dei polinomi \(p(x)\) e \(q(x)\), come illustrato di seguito:
\[p(x) = \sum_{i=0}^{n-1} a_i \, x^i = a_0 \, x^0 + a_1 \, x^1 + \dots + a_{n-1} \, x^{n-1} \qquad\quad q(x) = \sum_{i=0}^{m-1} b_i \, x^i = b_0 \, x^0 + b_1 \, x^1 + \dots + b_{m-1} \, x^{m-1}\]
\[p(x) \cdot q(x) = \sum_{i=0}^{n+m-2} c_i \, x^i \qquad\quad \text{con} \quad c_k = \!\! \sum_{\substack{i+j=k \\[2pt] 0 \leq i < n \\[2pt] 0 \leq j < m}} \!\! a_i b_j \qquad \text{per} \; k = 0, 1, \dotsc, n + m\]

\noindent
Applicando la pura definizione per il calcolo della convoluzione, la complessità risultante è quadratica nel degree bound dei
polinomi. La soluzione più efficiente che verrà descritta di seguito consiste in:
\begin{itemize}
	\item trasformare, tramite la DFT, i polinomi \(p(x)\) e \(q(x)\) dalla rappresentazione per vettori di coefficienti
	alla rappresentazione per vettori di punti o valori, valutando i polinomi nelle radici n-esime dell'unità, con complessità
	\(\Theta(n \cdot \log n)\) con \(n\) pari al degree bound dei polinomi
	\item eseguire il prodotto di Hadamard tra i vettori di punti o valori ottenuti, che equivale al prodotto di convoluzione
	dei vettori di coefficienti, con complessità lineare nel degree bound dei polinomi con complessità \(\Theta(n)\) con \(n\)
	pari al degree bound dei polinomi
	\item eseguire l'inversa della DFT per trasformare il vettore di punti o valori del polinomio risultante in un vettore
	di coefficienti con complessità \(\Theta(n \cdot \log n)\) con \(n\) pari al degree bound dei polinomi
\end{itemize}

\subsection{Rappresentazioni di un polinomio}
\subsubsection*{Rappresentazione tramite vettore di coefficienti}
Un polinomio \(p(x)\) di degree bound \(n\) può essere rappresentato tramite un vettore di coefficienti \(\underline{a}\)
di dimensione \(n\), in cui l'elemento \(a_i\) rappresenta il coefficiente del termine \(x^i\) del polinomio. Il prodotto
di due polinomi rappresentati in questo modo è dato dal prodotto di convoluzione tra i vettori di coefficienti.
\[\underline{a} = (a_0, a_1, \dotsc, a_{n-1}) \qquad \underline{b} = (b_0, b_1, \dotsc, b_{m-1}) \qquad \underline{c} = (c_0, c_1, \dotsc, c_{n+m-2})\qquad \underline{c} = \underline{a} * \underline{b}\]

\subsubsection*{Rappresentazione tramite vettore di punti o valori}
Un polinomio \(p(x)\) di degree bound \(n\) può essere rappresentato tramite un vettore di punti o valori \(\underline{y}\)
di dimensione \(n\), in cui l'elemento \(y_i\) rappresenta il valore del polinomio in un punto specifico \(x_i\). Esistono
infiniti modi di rappresentare un polinomio per punti, a seconda dei punti scelti, ma data una qualsiasi rappresentazione
esiste solo un polinomio di grado \(n\) associato a tale rappresentazione.

Il prodotto di due polinomi rappresentati per punti è dato dal prodotto element-wise tra i vettori di punti, avendo
l'accortezza di estenderli con nuovi elementi fino alla dimensione necessaria a rappresentare il polinomio risultante,
ovvero \(n + m - 1\) se i polinomi di partenza hanno degree bound \(n\) e \(m\).
\[(\underline{x}, \; \underline{y}) \text{ rappresentazione per punti di } p(x) \text{ con } \begin{array}{c}
	\underline{x} = (x_0, x_1, \dotsc, x_{n-1}) \\ \underline{y} = (y_0, y_1, \dotsc, y_{n-1})
\end{array} \text{ e } y_i = p(x_i)\]

\newpage

\subsection{Valutazione di un polinomio}
\subsubsection*{Valutazione di un polinomio usando la definizione}
Il processo di valutazione di un polinomio \(p(x)\) consiste nel calcolare il valore del polinomio in un punto specifico
\(x_i\). Tale processo permette di trasformare un polinomio dalla rappresentazione per vettore di coefficienti alla
rappresentazione per punti.

Il costo di valutazione di un polinomio \(p(x)\) di degree bound \(n\) in un punto \(x_i\) usando la definizione di polinomio
è lineare nel degree bound del polinomio.
\[T(n) \; = \underbrace{\; n - 2 \;}_{\text{potenze \(x^2, x^3, ...\)}} + \underbrace{\; n - 1 \;}_{\text{prodotti \(a_i \cdot x^i\)}} + \;\;\underbrace{\; n - 1 \;}_{\text{addizioni}} \; = \; 3n - 4 \; = \; \Theta(n)\]
Considerando che un polinomio deve essere valutato in \(n\) punti per essere trasformato in una rappresentazione per punti,
la complessità complessiva della trasformazione risulta quadratica.

\subsubsection*{Valutazione di un polinomio con l'algoritmo di Horner}
L'idea dell'algoritmo di Horner è quella di riscrivere un polinomio \(p(x)\) in una forma che permette di valutare
il polinomio ricorsivamente come segue:
\[p(x) = a_0 + a_1 x + a_2 x^2 + \dots + a_{n-1} x^{n-1} = a_0 + x \cdot \left( a_1 + a_2 x + \dots + a_{n-1} x^{n-2} \right)\]
L'equazione alle ricorrenze e la complessità dell'algoritmo di Horner per valutare il polinomio in un punto rimangono lineari
nel degree bound del polinomio, per cui non c'è un netto miglioramento rispetto all'approccio precedente.
\[T(n) \; = \; \begin{cases}
	\; 0 & \text{se } n = 1 \\[3pt]
	\; T(n-1) + 2 & \text{se } n > 1
\end{cases} \qquad\qquad T(n) = 2(n-1) = \Theta(n)\]

\subsubsection*{Valutazione di un polinomio sulle radici dell'unità}
Questo approccio si basa sulla possibilità di riutilizzare i conti fatti per valutare un polinomio \(p(x)\) in un punto
\(x_i\) per valutare lo stesso polinomio in altri punti \(x_j\), legati a \(x_i\) da una certa simmetria. Di seguito viene
illustrato un esempio per un polinomio di degree bound \(n = 8\).
\[\begin{array}{c c}
	p(x) \; = \; a_0 + a_1 x + a_2 x^2 + a_3 x^3 + a_4 x^4 + a_5 x^5 + a_6 x^6 + a_7 x^7 & \substack{\text{polinomio } p(x) \text{ di} \\ \text{degree bound } 8} \\[10pt]
	p(x) \; = \; \underbrace{(a_0 + a_2 x^2 + a_4 x^4 + a_6 x^6)}_{p^{[0]}(x^2)} \; + \; x \cdot \underbrace{(a_1 + a_3 x^2 + a_5 x^4 + a_7 x^6)}_{p^{[1]}(x^2)} & \substack{\text{si dividono le potenze} \\ \text{pari e dispari ottenendo} \\ \text{due polinomi in } x^2 \\ \text{di degree bound } 4} \\[20pt]
	p(x) = p^{[0]}(x^2) + x \cdot p^{[1]}(x^2) \qquad \begin{aligned} p(-x) &= p^{[0]}((-x)^2) - x \cdot p^{[1]}((-x)^2) \\ &= p^{[0]}(x^2) - x \cdot p^{[1]}(x^2) \end{aligned} & \substack{\text{le valutazioni del} \\\text{polinomio in } p(x) \text{ e} \\ p(-x) \text{ condividono gli} \\ \text{stessi due contribuiti} \\ p^{[0]}(x^2) \text{ e } p^{[1]}(x^2)} \\[22pt]
	\begin{array}{c}
		p^{[0]}(x^2) = p^{[0][0]}(x^4) + x^2 \cdot p^{[0][1]}(x^4) \\
		p^{[0]}(-x^2) = p^{[0][0]}(x^4) - x^2 \cdot p^{[0][1]}(x^4) \\[7pt]
		\begin{aligned}
			p(ix) &= p^{[0]}({ix}^2) + ix \cdot p^{[1]}({ix}^2) \\ &= p^{[0]}(-x^2) + ix \cdot p^{[1]}(-x^2)
		\end{aligned}
	\end{array} \begin{array}{c}
		p^{[1]}(x^2) = p^{[1][0]}(x^4) + x^2 \cdot p^{[1][1]}(x^4) \\
		p^{[1]}(-x^2) = p^{[1][0]}(x^4) - x^2 \cdot p^{[1][1]}(x^4) \\[7pt]
		\begin{aligned}
			p(ix) &= p^{[0]}({ix}^2) - ix \cdot p^{[1]}({ix}^2) \\ &= p^{[0]}(-x^2) - ix \cdot p^{[1]}(-x^2)
		\end{aligned}
	\end{array} & \substack{\text{calcolando i polinomi} \\ p^{[0]}(x^2) \text{ e } p^{[1]}(x^2) \\ \text{ si ottengono ``gratis''} \\ p^{[0]}(-x^2) \text{ e } p^{[1]}(-x^2) \\ \text{che forniscono anche} \\[1pt] \text{le due valutazioni} \\ p(ix) \text{ e } p(-ix) \\ \text{ con } i^2 = -1}
\end{array}\]
È possibile infine eseguire un'ulteriore passaggio ottenendo gli otto contributi \(p^{[i][j][k]}(x^8)\), ciascuno di
degree bound \(8/2/2/2 = 1\), che forniscono le valutazioni del polinomio \(p(x)\) anche per \(x \pm ix\) e \(-x \pm ix\).
Si valuta, \(p(x)\) sugli 8 punti \(\pm x, \; \pm ix,\; x \pm ix, \; -x \pm ix\) che sono proprio le radici ottave dell'unità.

Successivamente si dimostrerà che questo approccio (basato sulla DFT) ha complessità \(\Theta(n \cdot \log n)\) con \(n\)
pari al degree bound del polinomio, contro la complessità quadratica degli approcci precedenti.

\subsection{Interpolazione di un polinomio}
\subsubsection*{Interpolazione con la matrice di Vandermonde}
L'interpolazione di un polinomio \(p(x)\) consiste nel calcolare il vettore di coefficienti \(\underline{a}\) a partire
dal vettore di punti \(\underline{y}\). Tale processo permette di trasformare un polinomio dalla rappresentazione per
punti alla rappresentazione per vettore di coefficienti.

Un possibile approccio per l'interpolazione consiste nell'invertire la formula di valutazione di un polinomio di grado
\(n\), ottenendo così un sistema a \(n\) equazioni con \(n\) incognite (\(a_0, a_1, \dots a_{n-1}\)), i cui coefficienti
sono le potenze di \(x\).
\[p(x) = y = \sum_{i=0}^{n-1} a_i x^i \; \rightarrow \; \sum_{i=0}^{n-1} {x_j}^i a_i = y_j \; \rightarrow \; \underbrace{\begin{bmatrix}
	1 & x_0 & {x_0}^2 & \dots & {x_0}^{n-1} \\
	1 & x_1 & {x_1}^2 & \dots & {x_1}^{n-1} \\
	\vdots & \vdots & \vdots & \ddots & \vdots \\
	1 & x_{n-1} & {x_{n-1}}^2 & \dots & {x_{n-1}}^{n-1}
\end{bmatrix}}_{\text{matrice di Vandermonde} \;\; V(x_0, x_1, \dotsc, x_{n-1})} \cdot \begin{bmatrix}
	a_0 \\ a_1 \\ \vdots \\ a_{n-1}
\end{bmatrix} = \begin{bmatrix}
	y_0 \\ y_1 \\ \vdots \\ y_{n-1}
\end{bmatrix}\]
La matrice delle potenze di \(x\) è detta matrice di Vandermonde, denotata con \(V(x_0, x_1, \dots, x_{n-1})\). Il determinate
della matrice di Vandermonde è dato dalla seguente formula.
\[\det \left( V (\underline{x}) \right) \; = \prod_{0 \, \leq \, i \, < \, j \, \leq \, n-1} (x_j - x_i) \]
Se i punti \(x_0, x_1, \dotsc, x_{n-1}\) sono distinti, allora il determinante della matrice di Vandermonde è sempre
diverso da zero. Questo implica che la matrice è invertibile e il sistema di equazioni ha una soluzione unica. La complessità
rimane quadratica nel degree bound del polinomio dovuta al prodotto tra matrice inversa e vettore di punti, supponendo di
disporre dell'inversa della matrice di Vandermonde gratis.
\[\underline{a} = \left[ \; V(x_0, x_1, \dots x_{n-1}) \; \right]^{-1} \cdot \, \underline{y}\]

\subsubsection*{Interpolazione con il polinomio interpolatore di Lagrange}
Un altro approccio per l'interpolazione consiste nell'utilizzare il polinomio interpolatore di Lagrange
\[p(x) \; = \; \sum_{k=0}^{n-1} \; y_k \, L_k(x) \qquad\qquad L_i(x) \; = \; \frac{\displaystyle \prod_{j \neq i} (x - x_j)}{\displaystyle \prod_{j \neq i} (x_i - x_j)}\]
Si osserva che l'espressione in funzione di \(L_k(x)\) coincide con il polinomio \(p(x)\), in quanto, se tale polinomio
viene valutato in un punto \(x = x_k\) (con \(x_k \in\) rappresentazione per punti di \(p(x)\)), la sommatoria restituisce
proprio \(y_k\), che è proprio il valore del polinomio \(p(x)\) in \(x_k\), per ogni \(k \in [0, \, n-1]\), infatti:
\begin{itemize}
	\item \(L_i(x_k) = 1\) per \(i = k\), per la presenza di cancellazioni tra numeratore e denominatore
	\item \(L_i(x_k) = 0\) per ogni \(i \neq k\), siccome esisterà un fattore \(x_k - x_j = 0\) nel numeratore
	che non si semplificherà con nessun \(x_i - x_j\) del denominatore
\end{itemize}

\noindent
Il costo computazionale dell'approccio con il polinomio interpolatore di Lagrange rimane sempre quadratico nel degree bound
del polinomio da interpolare, in quanto è necessario calcolare \(n\) polinomi interpolatori di Lagrange \(L_k(x)\) (per
\(k \in [0, \, n-1]\)), ciascuno di grado \(n-1\).

\subsubsection*{Interpolazione di un polinomio sulle radici dell'unità tramite la matrice di Fourier}
Nei prossimi capitoli si osserverà che la valutazione di un polinomio \(p(x)\) sulle radici \(n\)-esime dell'unità
equivale al prodotto tra la matrice di Fourier e il vettore dei coefficienti del polinomio \(p(x)\). Di conseguenza,
per la trasformata inversa si utilizza l'inversa della matrice di Fourier e il vettore dei valori del polinomio \(p(x)\)
sulle radici \(n\)-esime dell'unità.

La valutazione (per come prima analizzato) costa \(\Theta(n \log_2 n)\), inoltre si dimostrerà che anche l'interpolazione costa \(\Theta(n \log_2 n)\),
ottenendo così un sistema efficiente per cambiare rappresentazione e utilizzare la rappresentazione per punti per il
calcolo del prodotto di polinomi.

\subsection{Matrice di Fourier e radici n-esime dell'unità}
\subsubsection*{Definizione di matrice di Fourier}
La matrice di Fourier \(F_n \in \mathbb{C}^{n \times n}\) è definita come la matrice di Vandermonde \(V\) i cui coefficienti
sono le radici \(n\)-esime dell'unità.
\[F_n = V\! \left( {\omega_n\!}^0, {\omega_n\!}^1, \dots, {\omega_n\!}^{n-1} \right) \qquad\qquad \text{con} \quad \omega_n = e^{-2\pi i / n} \;\; \text{radice principale dell'unità}\]

\subsubsection*{Radici n-esime dell'unità con relative proprietà e lemmi}
Una radice \(n\)-esima dell'unità è un numero complesso \(z_k = e^{i \varphi}\) con \(\varphi = 2\pi k/n\) e \(k \in [0, \; n\!-\!1]\):
\[{z_k}^n = [\exp (2 \pi i k / n)]^n = \exp (2 \pi i k / n \cdot n) = \exp (2 \pi i k) = 1\]
Una radice \(z_k\) è detta radice principale \(\omega_n\) se elevata per i vari valori di \(k \in [0, \; n\!-\!1]\)
genera tutte le restanti radici \(n\)-esime. Esistono più radici principali, solitamente si sceglie quella con \(k = 1\).
\[{\omega_n\!}^k = {z_1\!}^k = [\exp (2 \pi i / n)]^k = \exp (2 \pi i k / n) = z_k\]
Di seguito alcune proprietà e lemmi delle radici \(n\)-esime dell'unità:
\[\begin{array}{c c}
	{\omega_n}^i \cdot {\omega_n}^j = {\omega_n}^{i+j} & \qquad \begin{array}{c}\textbf{chiusura della moltiplicazione} \\ \text{l'insieme delle radici \(n\)-esime dell'unità} \\ \text{è chiuso rispetto alla moltiplicazione}\end{array} \\[15pt]
	{\omega_n}^k \cdot {\omega_n}^{n-k} = {\omega_n}^k \cdot {\omega_n}^{-k} = 1 & \qquad \begin{array}{c}\textbf{elemento inverso} \\ \text{l'insieme delle radici \(n\)-esime dell'unità possiede} \\ \text{l'elemento inverso } {\omega_n}^{n-k} \text{ rispetto alla moltiplicazione}\end{array} \\[15pt]
	{\omega_{d \cdot n}}^{d \cdot k} = {\omega_n}^k & \qquad \begin{array}{c}\textbf{lemma di cancellazione} \\ \text{moltiplicando per uno stesso fattore } d \text{ sia l'esponente} \\ \text{che il grado della radice, il risultato non cambia}\end{array} \\[15pt]
	{\omega_n}^{n/2} = {\omega_{2 \cdot n/2}}^{1 \cdot n/2} = \omega_2 & \qquad \begin{array}{c}\textbf{corollario del lemma di cancellazione} \end{array} \\[4pt]
	({\omega_n}^i)^2 = ({\omega_n}^{i + n/2})^2 = {\omega_{n/2}}^i & \qquad \begin{array}{c}\textbf{lemma di dimezzamento} \\ \text{le due radici } {\omega_n}^i \text{ e }{\omega_n}^{i + n/2} \text{ sono opposte} \\ \text{ed i loro quadrati sono uguali}\end{array} \\[15pt]
	\displaystyle \sum_{j=0}^{n-1} {({\omega_n}^i)}^j = \begin{cases} 0 & \text{se } i \; \text{mod} \; n = 0 \\[3pt] n & \text{altrimenti} \end{cases} & \qquad \begin{array}{c}\textbf{lemma di somma} \\ \text{la somma di tutte le potenze di una radice } n\text{-esima} \\ \text{dell'unità è zero se la radice è diversa da 1, altrimenti è \(n\)}\end{array}
\end{array}\]
\vspace{0pt}

\noindent
NOTA: il lemma di somma equivale alla somma degli elementi lungo la riga o la colonna \(i\)-esima della matrice di Fourier,
per cui la somma di tutte le righe e di tutte le colonne della matrice di Fourier è zero, tranne per la prima riga e la prima
colonna, la cui somma è \(n\).

\subsubsection*{Inversa della matrice di Fourier}
La matrice di Fourier è invertibile essendo una Vandermonde con determinante pari a \(n\) (per il lemma di somma) in quanto è
composta dalle radici \(n\)-esime dell'unità, tutte distinte tra loro.
\[F_{\omega_n}(i,j) = {\omega_n}^{ij} \qquad\qquad {F_{\omega_n}\!}^{-1}(i,j) = \frac{1}{n} \cdot {\omega_n}^{-ij} = \frac{1}{n} \cdot [F_{\omega_n}(i,j)]^{-1} = \frac{1}{n} \cdot F_{\omega_n}(n-i, \; j)\]
\vspace{-12pt}
\begin{align*}
	{F_n}^{-1} &= \frac{1}{n} \cdot V \!\left( \left({\omega_n}^{-1}\right)^0, \; \left({\omega_n}^{-1}\right)^1, \; \left({\omega_n}^{-1}\right)^2, \dots, \left({\omega_n}^{-1}\right)^{n-1} \right) \\
	&= \frac{1}{n} \cdot V \!\left( \left({\omega_n}^0\right)^{n-1}, \; \left({\omega_n}^1\right)^{n-1}, \; \left({\omega_n}^2\right)^{n-1}, \dots, \left({\omega_n}^{n-1}\right)^{n-1} \right) \\
	&= \frac{1}{n} \cdot V \!\left( {\omega_n}^0, \; {\omega_n}^{n-1}, \; {\omega_n}^{n-2}, \dots, {\omega_n}^1 \right)
\end{align*}

\noindent
La matrice inversa di Fourier è la matrice di Fourier scalata di \(1/n\) in cui le righe da \(1\) a \(n\!-\!1\) comprese
vengono invertite e l'insieme dei generatori è l'ultima riga della matrice di partenza.
\[({F_n\!}^{-\!1} \cdot F_n)(h, k) = \sum_{l=0}^{n-1} {F_n\!}^{-\!1}(h, l) \cdot F_n(l, k) = \sum_{l=0}^{n-1} \frac{1}{n} \, {\omega_n\!}^{-hl} \, {\omega_n\!}^{hk}= \frac{1}{n} \sum_{l=0}^{n-1} {\omega_n\!}^{l(k-h)} = \begin{cases} 1 & {\scriptstyle \text{se } h = k} \\ 0 & \text{\scriptsize altrimenti} \end{cases} = \; \mathbb{I}_n(h,k)\]

\subsection{Discrete Fourier Transform e algoritmo FFT}
\subsubsection*{Valutazione del polinomio sulle radici dell'unità con la matrice di Fourier}
Si osserva che il processo di valutazione di un polinomio \(p(x)\) sulle radici \(n\)-esime dell'unità, equivale al prodotto
tra la matrice di Fourier \(F_n\) e il vettore dei coefficienti \(\underline{a}\) del polinomio \(p(x)\). Come per la matrice
di Vandermonde, però, la complessità del prodotto applicando la pura definizione è \(\Theta(n^2)\).
\[\underline{y} \; = \begin{pmatrix}
	p(\omega_0) \\[3pt] p(\omega_1) \\[3pt] \vdots \\[3pt] p(\omega_{n-1})
\end{pmatrix} = \begin{pmatrix}
	\sum_{i=0}^{n-1} a_i \, {\omega_0}^i \\[3pt]
	\sum_{i=0}^{n-1} a_i \, {\omega_1}^i \\[3pt]
	\vdots \\[3pt]
	\sum_{i=0}^{n-1} a_i \, {\omega_{n-1}}^i
\end{pmatrix} =
\begin{pmatrix}
	{\omega_0}^0 & {\omega_0}^1 & \dots & {\omega_0}^{n-1} \\
	{\omega_1}^0 & {\omega_1}^1 & \dots & {\omega_1}^{n-1} \\
	\vdots & \vdots & \ddots & \vdots \\
	{\omega_{n-1}}^0 & {\omega_{n-1}}^1 & \dots & {\omega_{n-1}}^{n-1}
\end{pmatrix} \begin{pmatrix}
	a_0 \\ a_1 \\ \vdots \\ a_{n-1}
\end{pmatrix} = \; F_n \cdot \underline{a}\]

\subsubsection*{Formalizzazione della DFT}
Il prodotto tra la matrice di Fourier \(F_n\) e un generico vettore \(\underline{a}\) viene chiamato Discrete Fourier
Transform (DFT) del vettore \(\underline{a}\). È possibile, quindi, definire tale prodotto come una funzione come segue:
\[\underline{y} = F_n \cdot \underline{a} = \text{DFT}_n \left( \underline{a} \right)\]

\subsubsection*{Algoritmo efficiente per calcolare la DFT}
Si riprende la proprietà di sottostruttura individuata per valutare un polinomio \(p(x)\) per due generici punti \(x\) e
\(-x\) e la si applica per due generiche radici \(n\)-esime dell'unità opposte tra loro (\({\omega_n}^i, \; {\omega_n}^{i+n/2}\)).
\[\left\{ \; \begin{aligned}
	p(x) = p^{[0]}\!\left(x^2\right) + x \cdot p^{[1]}\!\left(x^2\right) \\
	p(-x) = p^{[0]}\!\left(x^2\right) - x \cdot p^{[1]}\!\left(x^2\right)
\end{aligned} \right. \quad \longrightarrow \quad \left\{ \; \begin{aligned}
	p \!\left({\omega_n}^i\right) = p^{[0]}\!\left({\omega_{n/2}}^i\right) + {\omega_n}^i \cdot p^{[1]}\!\left({\omega_{n/2}}^i\right) \\
	p \!\left({\omega_n}^{i + n/2}\right) = p^{[0]}\!\left({\omega_{n/2}}^i\right) - {\omega_n}^i \cdot p^{[1]}\!\left({\omega_{n/2}}^i\right)
\end{aligned} \right.\]
\[\text{con} \;\; \left({\omega_n}^i\right)^2 = \left({\omega_n}^{i + n/2}\right)^2 = {\omega_{n/2}}^i \;\; \text{per il lemma di dimezzamento o quadrato di radici opposte}\]

\noindent
Si esprime l'espressione sopra in funzione del vettore dei valori \(\underline{y}\) che assume il polinomio \(p(x)\) valutato
sulle radici \(n\)-esime dell'unità:
\[\underline{y} = \text{DFT}_n \left( \underline{a} \right) = \left\{ \;\; \begin{aligned}
	y_i = {y_i}^{[0]} + {\omega_n}^i \cdot {y_i}^{[1]} \\[4pt]
	y_{i + n/2} = {y_i}^{[0]} - {\omega_n}^i \cdot {y_i}^{[1]}
\end{aligned} \right. \quad \text{ per } i \in [0, \; n/2\!-\!1] \quad \text{con} \quad \begin{aligned}{}
	{y_i}^{[0]} &= p^{[0]} \! \left( {\omega_{n/2}}^i \right) \\ {y_i}^{[1]} &= p^{[1]} \! \left( {\omega_{n/2}}^i \right)
\end{aligned}\]

\noindent
Si osserva che i vettori \(\underline{y}^{[0]}\) e \(\underline{y}^{[1]}\) corrispondono alla valutazione dei polinomi
\(p^{[0]}(x)\) e \(p^{[1]}(x)\) sulle radici \(n/2\)-esime dell'unità, che possono essere calcolati riapplicando lo stesso
approccio in maniera ricorsiva.
\[y^{[0]} = F_{n/2} \cdot \underline{a}^{[0]} = \text{DFT}_{n/2} \! \left( \underline{a}^{[0]} \right) \qquad \text{e} \qquad y^{[1]} = F_{n/2} \cdot \underline{a}^{[1]} = \text{DFT}_{n/2} \! \left( \underline{a}^{[1]} \right)\]

\subsubsection*{Analisi della complessità dell'algoritmo efficiente per calcolare la DFT}
Il costo di ogni nodo interno è dato interamente dal costo del combine che comprende:
\begin{itemize}
	\item \(n/2\) moltiplicazioni tra i valori \(y_i^{[1]}\) e le radici \(n\)-esime dell'unità \({\omega_n}^i\)
	\item \(n/2\) addizioni per calcolare i valori \(y_i\)
	\item \(n/2\) sottrazioni per calcolare i valori \(y_{i + n/2}\)
	\item \(n/2 - 2\) prodotti per calcolare le radici \(n\)-esime dell'unità \({\omega_n}^i\) a partire da \(\omega_n\);
	il \(-2\) è dovuto al fatto che \({\omega_n}^0 = 1\) e \({\omega_n}^1 = \omega_n\) sono banali e non serve calcolarle
\end{itemize}
L'equazione alle ricorrenze e la complessità asintotica risultano, quindi, essere:
\[T(n) \; = \; \begin{cases}
	\; 0 & \text{se } n = 1 \\[3pt]
	\; 2 \cdot T(n/2) + 2n - 2 & \text{se } n > 1
\end{cases} \qquad\qquad T(n) = 2 n \log_2 n - 2 \log_2 n\]

\subsubsection*{Osservazione dei valori alle foglie dell'albero della ricorsione}
Costruendo l'albero della ricorsione si può osservare che i valori alle foglie sono ordinati secondo la permutazione
bit-reversal. Tale permutazione è ottenuta semplicemente invertendo l'ordine dei bit dell'indice \(i\) che identifica la foglia (es. \(i = 6\)
con rappresentazione binaria \(110\) diventa \(011\) che corrisponde a \(3\)).
\[\begin{array}{c c}
	\text{DFT}_{n}(\underline{a}) & (a_0, \; a_1, \; a_2, \; a_3, \; a_4, \; a_5, \; a_6, \; a_7) \\[3pt]
	\text{DFT}_{n/2}(\underline{a}^{[*]}) & (a_0, \; a_2, \; a_4, \; a_6) \quad\;\; (a_1, \; a_3, \; a_5, \; a_7) \\[3pt]
	\text{DFT}_{n/4}(\underline{a}^{[*][*]}) & (a_0, \, a_4) \quad (a_2, \, a_6) \quad\; (a_1, \, a_5) \quad (a_3, \, a_7) \\[3pt]
	\text{DFT}_{n/8}(\underline{a}^{[*][*][*]}) & (a_0) \;\; (a_4) \;\; (a_2) \;\; (a_6) \;\; (a_1) \;\; (a_5) \;\; (a_3) \quad (a_7)
\end{array}\]

\subsubsection*{Implementazione in pseudocodice dell'algoritmo FFT}
\begin{algo}{FFT}{$\underline{a}$}{rappresentazione tramite vettore di coefficienti $\underline{a}$ del polinomio $p(x)$}{rappresentazione tramite vettore di punti $\underline{y}$ del polinomio $p(x)$}
	\State \(n \gets \text{length}(\underline{a})\)
	\State
	\If{\(n = 1\)} \Comment{caso base}
		\State \textbf{return} \(a\)
	\EndIf

	\State
	\State \(\underline{a}^{[0]} \gets (a_0, a_2, a_4, \dots, a_{n-2})\) \Comment{fase di divide}
	\State \(\underline{a}^{[1]} \gets (a_1, a_3, a_5, \dots, a_{n-1})\)

	\State
	\State \(\underline{y}^{[0]} \gets \text{FFT}(\underline{a}^{[0]})\) \Comment{chiamate ricorsive}
	\State \(\underline{y}^{[1]} \gets \text{FFT}(\underline{a}^{[1]})\)

	\State
	\State \(\omega_n \gets \exp(2 \pi i /n)\) \Comment{fase di combine}
	\State \(\omega_i \gets 1\)
	\For{\(i \gets 0\)}{\(n/2-1\)}
		\State \(t \gets \omega_i \cdot y_i^{[1]}\)
		\State \(y_i \gets y_i^{[0]} + t\)
		\State \(y_{i + n/2} \gets y_i^{[0]} - t\)
		\State \(\omega_i \gets \omega_i \cdot \omega_n\)
	\EndFor
		
	\State \textbf{return} \(\underline{y}\)
\end{algo}

\newpage

\subsection{Moltiplicazione di polinomi con la DFT e padding dei coefficienti}
\subsubsection*{Padding dei coefficienti per il prodotto di polinomi}
Il prodotto di due polinomi \(p(x)\) e \(q(x)\) con degree bound rispettivamente di \(n\) e \(m\) è un polinomio
\(r(x) = p(x) \cdot q(x)\) con degree bound \(n + m - 1\). Per poter rappresentare correttamente tramite vettore di punti
il polinomio risultante è necessario valutare \(r(x)\) su almeno \(n + m - 1\) punti distinti.

Dal prodotto element-wise tra i vettori di punti di \(p(x)\) e \(q(x)\) valutati in \(\max\{n, m\}\) punti, si otterrà un
vettore di punti non sufficientemente lungo per rappresentare \(r(x)\). È necessario estendere le rappresentazioni tramite
vettori di coefficienti al degree bound \(n + m - 1\) del polinomio risultante, aggiungendo elementi nulli ai vettori di
coefficienti di \(p(x)\) e \(q(x)\). Questa operazione è detta padding dei coefficienti e gli elementi nulli aggiunti
corrispondo ai coefficienti delle potenze di grado superiore al degree bound dei due polinomi da moltiplicare.
\vspace{-5pt}
\begin{align*}
	\underline{a} = (a_0, a_1, \dots, a_{n-1}) \; \rightarrow \; \underline{a} = (a_0, a_1, \dots, a_{n-1}, \overbrace{0, 0, \dots, 0}^{m \; \text{elementi}}) \\
	\underline{b} = (b_0, b_1, \dots, b_{m-1}) \; \rightarrow \; \underline{b} = (b_0, b_1, \dots, b_{m-1}, \underbrace{0, 0, \dots, 0}_{n \; \text{elementi}}) \\
\end{align*}
\vspace{-33pt}

\subsubsection*{Padding per raggiungere la potenza di due più vicina}
Siccome l'algoritmo FFT è stato strutturato per funzionare con polinomi di degree bound pari ad una potenza di due,
è necessario estendere i vettori di coefficienti al degree bound pari alla potenza di due maggiore più vicina al degree
bound del polinomio risultante. Sarà necessario, quindi, effettuare un ulteriore padding e aggiungere ulteriori elementi
nulli ai vettori di coefficienti.
\begin{align*}
	\underline{a} = (a_0, a_1, \dots, a_{n-1}, 0, 0, \dots, 0) \; \rightarrow \; \underline{a} = \overbrace{(a_0, a_1, \dots, a_{n-1}, 0, 0, \dots, 0)}^{2^k \, \geq \, n + m - 1 \; \text{elementi}}) \\
	\underline{b} = (b_0, b_1, \dots, b_{m-1}, 0, 0, \dots, 0) \; \rightarrow \; \underline{b} = \underbrace{(b_0, b_1, \dots, b_{m-1}, 0, 0, \dots, 0)}_{2^k \, \geq \, n + m - 1 \; \text{elementi}}) \\
\end{align*}
\vspace{-33pt}

\subsubsection*{Note sulla complessità dopo il padding}
Definita \(N = 2^k\) la dimensione dei vettori di coefficienti dopo i due padding, la complessità del prodotto di polinomi
tramite la DFT è \(\Theta(N \log_2 N)\).

Supponendo che i degree bound dei polinomi \(p(x)\) e \(q(x)\) siano entrambi \(n\), il primo padding porterà la
dimensione dei vettori a \(2n\), pari al degree bound del polinomio risultante. La potenza più vicina a \(2n\) potrà
essere al massimo \(2 \cdot 2n = 4n\), siccome due potenze di due consecutive differiscono di un fattore 2.

Di conseguenza, la dimensione finale sarà al massimo \(N = 4n\), che provocherà un peggioramento di circa quattro volte
della complessità del prodotto, ma senza alterare quella asintotica di \(\Theta(n \log_2 n)\).

\subsubsection*{Note sull'interpretazione del padding in generale}
In generale non sempre è possibile applicare un padding all'input per ottenere una taglia di dimensione ottimale. Di seguito
alcune osservazioni utili per interpretare il padding in generale:
\begin{itemize}
	\item nel prodotto tra polinomi, il padding di elementi nulli non altera il polinomio, in quanto equivale ad aggiungere
	coefficienti nulli delle potenze di grado superiore
	\item nel mergesort, il padding potrebbe essere fatto con valori facili da discriminare (es. \(-\infty\) o \(+\infty\))
	per cui è possibile ottenere la stringa ordinata correttamente dopo averli eliminati dall'output
	\item nella trasformata di Fourier (e nella DCT), in generale, gli elementi da aggiungere al vettore di coefficienti
	affinché l'output non venga alterato sono molto complessi da determinare, per cui non è possibile applicare un padding
	in maniera efficiente
\end{itemize}

\noindent
NOTA: il padding di zeri nei vettori di input della DFT per la moltiplicazione di polinomi altera il risultato della DFT
perché banalmente si sta valutando il polinomio su più radici dell'unità. Tuttavia, in questo contesto non è un problema,
in quanto entrambe le rappresentazioni ottenute prima e dopo il padding (seppur sostanzialmente diverse) sono equivalenti
per rappresentare lo stesso polinomio di partenza.

\subsubsection*{Padding per \(k\)-potenza di un polinomio}
Si osserva che man mano che il polinomio viene elevato a potenza, il degree bound continua a crescere e sarebbe necessario
effettuare continui padding per garantire la valutazione su un numero di punti sufficientemente grande per rappresentare
i risultati intermedi e tale numero deve anche essere pari ad una potenza di 2.

In alternativa a tale processo di padding continuo, si può effettuare un padding iniziale già sufficientemente grande
per garantire che tutti i risultati intermedi e finali siano rappresentati correttamente.

Analizzando la complessità di entrambi i processi, si osserva che il costo complessivo dell'elevamento a potenza è
praticamente lo stesso ed è pari a \(\Theta(2 \, kn \log_2 kn)\), che non altera la complessità asintotica.

\subsection{Osservazioni utili per l'esame}
\subsubsection*{Vettore di output di una trasformazione data la matrice e un certo vettore di input}
Quando viene richiesto di calcolare il vettore di output di una trasformazione data la matrice relativa e un certo vettore
di input, è utile osservare che il vettore di output è dato dalla combinazione lineare delle colonne della matrice pesate
dai coefficienti del vettore di input.
\[\underline{y} = F \cdot \underline{a} = \sum_{i=0}^{n-1} F(*,i) \cdot a_i \qquad \text{con } F(*,i) \text{ colonna } i\text{-esima di } F\]
Nel caso della trasformata di Fourier esistono alcuni vettori notevoli:
\begin{itemize}
	\item \(\underline{a} = [1, 0, 0, \dots 0] \;\; \rightarrow \;\; \underline{y} = F_n \cdot \underline{a} = [1, 1, 1, \dots, 1]\)
	siccome la prima colonna della matrice di Fourier è composta da tutti 1
	\item \(\underline{a} = [0, 1, 0, 0, \dots 0] \;\; \rightarrow \;\; \underline{y} = F_n \cdot \underline{a} = [{\omega_n}^0, {\omega_n}^1, {\omega_n}^2, \dots, {\omega_n}^{n-1}]\)
	siccome la seconda colonna della matrice di Fourier è composta dalle radici \(n\)-esime dell'unità
	\item \(\underline{a} = [1, 1, 1, \dots 1] \;\; \rightarrow \;\; \underline{y} = F_n \cdot \underline{a} = [n, 0, 0, \dots, 0]\)
	siccome la somma degli elementi lungo le righe della matrice di Fourier è zero, tranne per la prima riga, la cui somma è \(n\) (lemma di somma)
	\item \(\underline{a} = (n,k)\text{-sparso}\) con \(n\) e \(k\) potenze di 2, è un vettore di \(n\) elementi con \(x_i = 0 \text{ per } i \text{ mod } k \neq 0\).
	Un caso che compare spesso sono i vettori \((8,4)\text{-sparsi} : \underline{a} = [p, 0, 0, 0, q, 0, 0, 0]\) la cui trasformata è:
	\(\underline{y} = F_8 \cdot [p, 0, 0, 0, q, 0, 0, 0] = p \cdot [1,1,1,1,1,1,1,1] + q \cdot [1,-1,1,-1,1,-1,1,-1]\)
\end{itemize}

\subsubsection*{Composizione di trasformate di Fourier}
La composizione di due trasformate di Fourier \({F_n}^2\) provoca uno scalamento di un fattore \(n\) e un ribaltamento degli
elementi del vettore di input, escluso il primo elemento, che rimane invariato.
\[{F_n}^2 \cdot \underline{a} = n \cdot (a_0, a_{n-1}, a_{n-2}, \dots, a_1) = n \cdot \underline{a}^*\]

\noindent
Siccome il ribaltamento applicato due volte riporta gli elementi del vettore di input alla posizione originale, la composizione
di quattro trasformate di Fourier \({F_n}^4\) provoca soltanto uno scalamento di \(n^2\) senza modificare la posizione degli
elementi del vettore di input.
\[{F_n}^4 \cdot \underline{a} = n^2 \cdot (a_0, a_1, a_2, \dots, a_{n-1}) = n^2 \cdot \underline{a}\]

\subsubsection*{Inversa della trasformata di Fourier}
Dalla composizione di due trasformate di Fourier si ottiene che l'inversa della trasformata di Fourier di un vettore
\(\underline{y}\) è data dalla trasformata di Fourier del vettore \(\underline{y}\) scalato di \(1/n\), a cui vengono
ribaltati gli elementi, escluso il primo.
%\[\begin{array}{c c c}
%	1. & \quad \underline{y} = F_n \cdot \underline{a} \;\; \rightarrow \;\; \underline{y} = \text{DFT}_n(\underline{a}) &
%	\quad \substack{\text{si riscrive la trasformata} \\ \text{come funzione matematica}} \\[7pt]
%	2. & \quad n \cdot \underline{a}^* = {F_n}^2 \cdot \underline{a} \;\; \rightarrow \;\; n \cdot \underline{a}^* = \text{DFT}_n(\text{DFT}_n(\underline{a})) &
%	\quad \substack{\text{dalla formula di composizione} \\ \text{di due trasformate di Fourier}} \\[7pt]
%	3. & \quad n \cdot \underline{a}^* = \text{DFT}_n(\text{DFT}_n(\underline{a})) = \text{DFT}_n(\underline{y}) &
%	\quad \substack{\text{dalla formula 1.}\; : \; \text{DFT}_n(\underline{a}) \; = \; \underline{y}} \\[7pt]
%	4. & \quad \underline{a}^* = \frac{1}{n} \cdot \text{DFT}_n(\underline{y}) &
%	\quad \substack{\text{invertendo la formula sopra}}
%\end{array}\]
\[n \cdot \underline{a}^* = {F_n}^2 \cdot \underline{a} = \text{DFT}_n(\text{DFT}_n(\underline{a})) = \text{DFT}_n(\underline{y}) \implies \underline{a}^* = \frac{1}{n} \cdot \text{DFT}_n(\underline{y}) \]
\[\text{con} \quad \underline{y} = \text{DFT}_n(\underline{a}) \quad \text{e} \quad \underline{a}^* = \left( a_0, a_{n-1}, a_{n-2}, \dots a_1 \right)\]
